% Generated draft: requires mathematical and source review.
\documentclass[11pt]{article}
\usepackage[margin=25mm]{geometry}
\usepackage{fontspec}
\setmainfont{Latin Modern Roman}
\usepackage{amsmath,amssymb,mathtools,unicode-math}
\setmathfont{Latin Modern Math}
\usepackage{enumitem,needspace,xurl,hyperref}
\setlength{\parindent}{0pt}
\setlength{\parskip}{5pt}
\setlength{\emergencystretch}{3em}
\newcommand{\N}{\mathbb{N}}
\newcommand{\Z}{\mathbb{Z}}
\newcommand{\Q}{\mathbb{Q}}
\newcommand{\R}{\mathbb{R}}
\newcommand{\C}{\mathbb{C}}
\newcommand{\F}{\mathbb{F}}
\newcommand{\A}{\mathbb{A}}
\newcommand{\PP}{\mathbb{P}}
\newcommand{\E}{\mathbb{E}}
\newcommand{\eps}{\varepsilon}
\newcommand{\dd}{\,\mathrm{d}}
\begin{document}
\section*{UnsolvedMath ID: 800461}
\textbf{Source problem code:} catalogue entry
\subsection{Definitions and mathematical statement}
Let \((M_i^n,g_i,p_i)\) be a sequence of pointed \(n\)-dimensional Riemannian manifolds with limit \((X,d,p)\), as in the supplied source. The usual interpretation of the displayed convergence is pointed Gromov--Hausdorff convergence. Suppose that, uniformly in \(i\),
\[
\operatorname{Ric}_{g_i}\geq -\lambda g_i,
\qquad
\operatorname{Vol}_{g_i}(B_1(p_i))>v>0,
\]
where \(\lambda\) and \(v\) are fixed constants and \(B_1(p_i)\) is the open metric ball of radius \(1\).

A tangent cone at \(x\in X\) means any pointed Gromov--Hausdorff limit
\[
(X,r_j^{-1}d,x)\longrightarrow (T,d_T,o),
\qquad r_j>0,\quad r_j\to0.
\]
Under the customary convention, a pointed metric space is \(m\)-symmetric if it is pointed isometric to
\[
\bigl(\R^m\times C(Z),d_{\mathrm{prod}},(0,o_C)\bigr),
\]
where \(C(Z)\) is a metric cone with vertex \(o_C\), and \(d_{\mathrm{prod}}\) is the Euclidean product metric. The residual cone may be a single point. Here ``at least \(m\)-symmetric'' means that such a splitting with an \(\R^m\) factor exists.

\textbf{Conjecture 5 (attributed to Naber, 2020).} For every integer \(k\) with \(0\leq k\leq n-2\), there exists a subset \(\widehat{\mathcal S}^{\,k}\subseteq X\) such that
\[
\dim_{\mathrm H}\widehat{\mathcal S}^{\,k}\leq k
\]
and, for every \(x\in X\setminus\widehat{\mathcal S}^{\,k}\), every tangent cone at \(x\) is at least \((k+1)\)-symmetric. The residual cone and the splitting may depend on the tangent cone.

\textit{Unresolved source qualifications.} The supplied excerpt does not explicitly specify the convergence topology, the meaning of dimension, or the definition of symmetry. The statement above uses the customary pointed Gromov--Hausdorff, Hausdorff dimension, and metric cone splitting conventions, which require confirmation against the original source. The excerpt writes only \(k\leq n-2\); the nonnegative integer range is the conventional interpretation. Completeness and boundary assumptions on the manifolds, and the permitted range of \(\lambda\), are not supplied. No additional such hypotheses are asserted here.
\subsection{Short English statement}
The conjecture predicts that, for each indicated \(k\), an exceptional set of Hausdorff dimension at most \(k\) contains every point having any tangent cone without a Euclidean factor of dimension \(k+1\). It concerns every tangent cone at every point outside that set. The geometric conventions and missing manifold hypotheses require source review.
\subsection{Sources}
\begin{itemize}
\item Local UnsolvedMath registry, permanent ID \(800461\), exact source problem code \texttt{catalogue entry}, title ``Spaces with Lower Ricci Curvature --- Conjecture 5.'' The supplied registry attributes the conjecture to Naber and gives the proposed year as 2020.
\item Naber, \emph{Conjectures and Open Questions on Spaces with Lower Ricci Curvature} (2020), Conjecture 5. Supplied source URL: \url{https://arxiv.org/abs/2010.10031}. The background records access on 2026-10-16 and extraction by \texttt{arxiv\_latex}; these are supplied provenance statements, not independently verified access or source checks.
\item The supplied registry labels this problem ``open'' and records upstream revision \texttt{d371ee917ac9647fbb06cd6fc529deec6f3dda82}. The status is supplied curation, not an independently verified mathematical or literature assessment.
\end{itemize}
\end{document}
